% Cote 119: offprints, reports, courses, CV and letters (typed and printed part): archive cote 119, site page 204 (PDF page 205).
% typescript. Draft compared: none.
% AI-transcribed by Claude Opus 5.5 (claude-opus-5-5, adaptive thinking, effort high), Message Batches API, batch r119_full (msgbatch_01DwZNm5Y7YsKBWyAN4sFs8w), 2026-11-13; not yet checked by eye.
\scanpage{119/204}
\mspage{- 9 -}

que O. Zariski, J.P. Serre, H. Hironaka, D. Mumford, I. Manin, F. Chafarévitch.

b) \uline{Théorie locale des schémas et des morphismes de schémas}: Dans ce contexte
se placent les développements d'algèbre commutative mentionnés au par. 4, et
l'étude détaillée de notions comme celles de morphisme lisse, étale, net, plat ,
etc. Les quatre volumes de [15, Chap. IV] sont consacrés à ces développements,
qui ont d'ailleurs inspiré des développements analogues en théorie des espaces
analytiques et rigide-analytiques.

c) \uline{Techniques de construction de schémas}: Parmi les techniques développées,
exposées surtout dans [15 bis] et des séminaires non publiés (par moi-même et
d'autres), il y a la \uline{théorie de la descente} (cf. aussi [16, SGA I, Exp. V, VI]),
celle des \uline{schémas quotients}, des \uline{schémas de Hilbert}, des \uline{schémas de Picard},
des \uline{``modules'' formels}, le \uline{théorème d'existence} des faisceaux de modules algébri-%
ques associés à des modules formels ([15, Chap. III, par. 5]). Le point de vue
adopté est surtout celui de la construction d'un schéma à partir du foncteur
qu'il représente. Dans cette optique, je n'étais pas parvenu à une véritable
caractérisation maniable des foncteurs représentables par un schéma relatif
(localement de type fini sur un schéma noethérien) - c'est M. Artin qui y est
parvenu ultérieurement [*], en remplaçant la notion de schéma par celle, plus
générale et plus stable, d'espace algébrique. Parmi d'autres recherches dans
la même direction, suscitées par mes travaux, il y a celles de P. Murre sur les
schémas de Picard sur un corps [*],celles de D. Mumford et de M. Raynaud [*]
sur ces mêmes schémas sur des bases générales, et dans une certaine mesure
ceux de D. Mumford [*] et de S. Seshadri sur le passage au quotient, pour n'en
citer que quelques-uns.

d) \uline{Théories cohomologiques}:

1°) \uline{Cohomologie ``cohérente''} : résultats de finitude, de comparaison avec la
cohomologie formelle , cf. [15, Chap. III]. Théorèmes de dualité et des résidus :
un exposé systématique de mes idées et résultats est développé dans le séminaire
de R. Hartshorne [*], cf. aussi [18 bis].

2°) \uline{Cohomologie $\ell$-adique} : définition de la cohomologie étale, théorèmes de
comparaison, de finitude, de dimension cohomologique, de Lefschetz faible,
[16 SGA 4]; théorèmes de dualité, formules de Lefschetz et d'Euler-Poincaré,
application aux fonction L , [16, SGA 15].

3°) \uline{Cohomologie de De Rham} : [16 bis], [17 bis].

\margin{204}
